\newif\ifglobal

%%%%%%%%%%%%%%%%%%%%%%%
%%% Compile to view %%%
%%%%%%%%%%%%%%%%%%%%%%%

%%%%%%%%%%%%%%%%%%%%%%%%%%%%%%%%%%%%%%%%%%%%%%%%%%%
%%% Readme:                                     %%%
%%% https://www.overleaf.com/read/bwdjvfjksvjy/ %%%
%%%%%%%%%%%%%%%%%%%%%%%%%%%%%%%%%%%%%%%%%%%%%%%%%%%

% >>> M?: marks a module setting number or important notice

% >>> M4: Set {./relative-path} to external files for this module <<<
\def \modulefiles {./07-EFUSE}
% To add an external file, use:
% (Replace '---' with actual filename)
% '\input{\modulefiles/tables/---.tex}' for tables
% '\includegraphics{\modulefiles/figures/---}' for figures
% '\subfile{\modulefiles/---}' for register files

% >>> M1: This module is in global project? <<<
% 'YES' - uncomment; 'NO' - comment out
\globaltrue

\ifglobal
    \documentclass[main\_\_EN.tex]{subfiles}

\else
    \documentclass[openany, 10pt]{book}

    % >>> M2: In ./00-shared/preamble-trm-module, also find
    %         settings for visibility of todo notes, watermark <<<
    \usepackage{preamble-trm-module}

    % >>> M3: Temporary labels in English,
    %         you can add your own labels there <<<
    \myexternaldocument{./00-shared/temp-labels-en}

\fi %\ifglobal


\setcounter{secnumdepth}{4}

% >>> M5: If this module needs any updates in
% './00-shared/preamble-trm-module.sty'
% (include additional package, etc.), contact your module PM
% or create the file './00-shared/changelog.tex' make the
% necessary updates yourself and document them in this file <<<


\begin{document}


% Sets language into which pre-defined bits of text are translated
\selectlanguage{english}

% Do not delete this block
\ifglobal
% Add the following front matter if
% this module is outside of global project
\else
    \InputIfFileExists{readme}{}{}
    {\let\clearpage\relax \listoftodos}
    \tableofcontents
    %\listoftables
    %\listoffigures
\fi


%%%%%%%%%%%%%%%%%%%%%%%%%
%%% TRM Module Begins %%%
%%%%%%%%%%%%%%%%%%%%%%%%%

\hypertarget{efuse}{}
\chapter{eFuse Controller (EFUSE)}
\label{mod:efuse}

\section{Overview}

The \chipname{} contains a 4-Kbit eFuse to store parameters. These parameters are burned and read by an eFuse Controller. Once an eFuse bit is programmed to 1, it can never be reverted to 0. The eFuse controller programs individual bits of parameters in eFuse according to users configurations. Some of these parameters can be read by users using the eFuse controller. If read-protection for some data is not enabled, that data is readable from outside the chip. If read-protection is enabled, that data can not be read from outside the chip. In all cases, however, some keys stored in eFuse can still be used internally by hardware cryptography modules such as Digital Signature, HMAC, etc., without exposing this data to the outside world.

\section{Features}

\begin{itemize}
    \item 4-Kbit in total, with 1792 bits available for users
    \item One-time programmable storage
    \item Configurable write protection
    \item Configurable read protection
    \item Various hardware encoding schemes protect against data corruption
\end{itemize}

\section{Functional Description}

\subsection{Structure}

The eFuse data is organized in 11 blocks (BLOCK0 \verb+~+ BLOCK10). BLOCK0 has 640 bits in totall. BLOCK1 has 288 bits and each block of BLOCK2 \verb+~+ 10 has 352 bits.

BLOCK0, which holds most parameters, has 25 bits that are readable but useless to users (the details are showed in Section \ref{sec:block0_hw_bits}), and 29 further bits are reserved for future use.

Table \ref{tab:efuse-block0-para} lists all the parameters in BLOCK0 and their offsets, bit widths, as well as information on whether they can be used by hardware, which bits are write-protected, and corresponding descriptions.

The {\bfseries \hyperref[fielddesc:EFUSEWRDIS]{EFUSE\_WR\_DIS}} parameter is used to disable the writing of other parameters, while {\bfseries \hyperref[fielddesc:EFUSERDDIS]{EFUSE\_RD\_DIS}} is used to disable users from reading BLOCK4 \verb+~+ BLOCK10. For more information on these two parameters, please see Section \ref{sec:efuse-wr-dis} and Section \ref{sec:efuse-rd-dis}.


\clearpage
\begin{landscape}

\begin{longtable}[l]{ |m{8.2cm}|r|c|r|m{9cm}| }    
\caption{Parameters in eFuse BLOCK0}
    \label{tab:efuse-block0-para} \\\hline
    
    \rowcolor{lightgray}
    \textbf{Parameters} & \textbf{\thead{Bit \\Width}} & \textbf{\thead{Accessible \\by Hardware}} & \textbf{\thead{Programming-Protection \\by \hyperref[fielddesc:EFUSEWRDIS]{EFUSE\_WR\_DIS}\\ Bit Number}}
    & \textbf{Description} \\
    \endfirsthead
    
    \multicolumn{5}{c}{\bfseries \tablename\ \thetable{} -- cont'd from previous page}\\\hline
    \rowcolor{lightgray}
    \textbf{Parameters} & \textbf{\thead{Bit \\Width}} & \textbf{\thead{Accessible \\by Hardware}} & \textbf{\thead{Programming-Protection \\by \hyperref[fielddesc:EFUSEWRDIS]{EFUSE\_WR\_DIS}\\ Bit Number}}
    & \textbf{Description} \\
    \endhead 
    \multicolumn{5}{r}{\textbf{Cont'd on next page}} \\
    \endfoot
    \endlastfoot
    \hline   
    \hyperref[fielddesc:EFUSEWRDIS]{\textbf{EFUSE\_WR\_DIS}} & 32 & Y & N/A & Represents whether writing of individual eFuses is disabled. \\
    \hline
    \hyperref[fielddesc:EFUSERDDIS]{\textbf{EFUSE\_RD\_DIS}} & 7 & Y & 0  & Represents whether users' reading from BLOCK4 \verb+~+ 10 is disabled. \\
    \hline
    \hyperref[fielddesc:EFUSEDISICACHE]{EFUSE\_DIS\_ICACHE} &  1 & Y & 2  & Represents whether iCache is disabled. \\
    \hline
    \hyperref[fielddesc:EFUSEDISDCACHE]{EFUSE\_DIS\_DCACHE} & 1 & Y & 2  & Represents whether dCache is disabled. \\
    \hline
    \hyperref[fielddesc:EFUSEDISDOWNLOADICACHE]{EFUSE\_DIS\_DOWNLOAD\_ICACHE} & 1 & Y & 2 & Represents whether iCache is disabled in Download mode. \\   
    \hline
    \hyperref[fielddesc:EFUSEDISDOWNLOADDCACHE]{EFUSE\_DIS\_DOWNLOAD\_DCACHE} & 1 & Y & 2 & Represents whether dCache is disabled in Download mode. \\      
    \hline    
    \hyperref[fielddesc:EFUSEDISFORCEDOWNLOAD]{EFUSE\_DIS\_FORCE\_DOWNLOAD} & 1 & Y & 2 & Represents whether the function that forces chip into download mode is disabled. \\   
    \hline
    \hyperref[fielddesc:EFUSEDISUSBOTG]{EFUSE\_DIS\_USB\_OTG} &1 & Y & 2 & Represents whether USB OTG function is disabled. \\     
    \hline
    \hyperref[fielddesc:EFUSEDISTWAI]{EFUSE\_DIS\_TWAI} & 1 & Y & 2 & Represents whether TWAI Controller is disable. \\  
    \hline
    \hyperref[fielddesc:EFUSEDISAPPCPU]{EFUSE\_DIS\_APP\_CPU} & 1 & Y & 2 & Represents whether app CPU is disable. \\   
    \hline
    \hyperref[fielddesc:EFUSESOFTDISJTAG]{EFUSE\_SOFT\_DIS\_JTAG} & 3 & Y & 31 & Represents whether JTAG with soft-disable is disabled. \\
    \hline
    \hyperref[fielddesc:EFUSEDISPADJTAG]{EFUSE\_DIS\_PAD\_JTAG} & 1 & Y & 2 & Represents whether pad JTAG is permanently disabled. \\   
    \hline
    \hyperref[fielddesc:EFUSEDISDOWNLOADMANUALENCRYPT]{EFUSE\_DIS\_DOWNLOAD\_
    MANUAL\_ENCRYPT} &  1 & Y & 2 & Represents whether flash encryption is disabled in Download boot modes.\\ 
    \hline
    \hyperref[fielddesc:EFUSEUSBEXCHGPINS]{EFUSE\_USB\_EXCHG\_PINS} & 1 & Y & 30 & Represents whether USB D+ and D- pins are swapped. \\
    \hline
    \hyperref[fielddesc:EFUSEEXTPHYENABLE]{EFUSE\_EXT\_PHY\_ENABLE} & 1 & N & 30 & Represents whether external USB PHY is disabled. \\
    \hline
    \hyperref[fielddesc:EFUSEVDDSPIXPD]{EFUSE\_VDD\_SPI\_XPD} & 1 & Y & 3 & Represents whether Flash Voltage Regulator is powered up. \\        
    \hline
    \hyperref[fielddesc:EFUSEVDDSPITIEH]{EFUSE\_VDD\_SPI\_TIEH} & 1 & Y & 3 & Represents whether Flash Voltage Regulator output is short connected to VDD3P3\_RTC\_IO. \\   
    \hline
    \hyperref[fielddesc:EFUSEVDDSPIFORCE]{EFUSE\_VDD\_SPI\_FORCE} & 1 & Y & 3 & Represents whether to force using EFUSE\_VDD\_SPI\_XPD and EFUSE\_VDD\_SPI\_TIEH to configure flash voltage LDO. \\ 
    \hline
    \hyperref[fielddesc:EFUSEWDTDELAYSEL]{EFUSE\_WDT\_DELAY\_SEL} & 2 & Y & 3 & Represents RTC watchdog timeout threshold. \\
    \hline
    \hyperref[fielddesc:EFUSESPIBOOTCRYPTCNT]{EFUSE\_SPI\_BOOT\_CRYPT\_CNT} & 3 & Y & 4 & Represents whether SPI boot encrypt/decrypt is disabled. \\
    \hline
    \hyperref[fielddesc:EFUSESECUREBOOTKEYREVOKE0]{EFUSE\_SECURE\_BOOT\_KEY\_
    REVOKE0} & 1 & N & 5 & Represents whether the first secure boot key is revoked. \\
    \hline
    \hyperref[fielddesc:EFUSESECUREBOOTKEYREVOKE1]{EFUSE\_SECURE\_BOOT\_KEY\_
    REVOKE1} & 1 & N & 6 & Represents whether the second secure boot key is revoked. \\
    \hline
    \hyperref[fielddesc:EFUSESECUREBOOTKEYREVOKE2]{EFUSE\_SECURE\_BOOT\_KEY\_
    REVOKE2} & 1 & N & 7 & Represents whether the third secure boot key is revoked. \\
    \hline
    \hyperref[fielddesc:EFUSEKEYPURPOSE0]{EFUSE\_KEY\_PURPOSE\_0} & 4 & Y & 8 & Represents Key0 purpose, see Table \ref{tab:efuse-key-purpose}. \\
    \hline
    \hyperref[fielddesc:EFUSEKEYPURPOSE1]{EFUSE\_KEY\_PURPOSE\_1} & 4 & Y & 9 & Represents Key1 purpose, see Table \ref{tab:efuse-key-purpose}.\\
    \hline
    \hyperref[fielddesc:EFUSEKEYPURPOSE2]{EFUSE\_KEY\_PURPOSE\_2} & 4 & Y & 10 & Represents Key2 purpose, see Table  \ref{tab:efuse-key-purpose}. \\ 
    \hline
    \hyperref[fielddesc:EFUSEKEYPURPOSE3]{EFUSE\_KEY\_PURPOSE\_3} & 4 & Y & 11 & Represents Key3 purpose, see Table \ref{tab:efuse-key-purpose}. \\    
    \hline
    \hyperref[fielddesc:EFUSEKEYPURPOSE4]{EFUSE\_KEY\_PURPOSE\_4} & 4 & Y & 12 & Represents Key4 purpose, see Table \ref{tab:efuse-key-purpose}.\\    
    \hline
    \hyperref[fielddesc:EFUSEKEYPURPOSE5]{EFUSE\_KEY\_PURPOSE\_5} & 4 & Y & 13 & Represents Key5 purpose, see Table \ref{tab:efuse-key-purpose}. \\  
    \hline
    \hyperref[fielddesc:EFUSESECUREBOOTEN]{EFUSE\_SECURE\_BOOT\_EN} & 1 & N & 15 & Represents whether secure boot is enabled. \\ 
    \hline
    \hyperref[fielddesc:EFUSESECUREBOOTAGGRESSIVEREVOKE]{EFUSE\_SECURE\_BOOT\_AGG
    RESSIVE\_REVOKE} & 1 & N & 16 & Represents whether aggressive revoke of secure boot keys is enabled. \\
    \hline
    \hyperref[fielddesc:EFUSEDISUSBJTAG]{EFUSE\_DIS\_USB\_JTAG} & 1 & Y & 2 & Represents whether the function of usb\_serial\_jtag that switch usb to jtag is disabled. \\
    \hline
    \hyperref[fielddesc:EFUSEDISUSBSERIALJTAG]{EFUSE\_DIS\_USB\_SERIAL\_JTAG} & 1 & Y & 2 & Represents whether usb\_serial\_jtag function is disabled. \\
    \hline
    \hyperref[fielddesc:EFUSESTRAPJTAGSEL]{EFUSE\_STRAP\_JTAG\_SEL} & 1 & Y & 2 & Represents whether to enable selection between usb\_to\_jtag or pad\_to\_jtag through GPIO3. 0: pad\_to\_jtag; 1: usb\_to\_jtag. \\
    \hline
    \hyperref[fielddesc:EFUSEUSBPHYSEL]{EFUSE\_USB\_PHY\_SEL} & 1 & Y & 2 & Represents the connection relationship between internal PHY, external PHY, and USB OTG, USB Serial/JTAG. \\
    \hline
    \hyperref[fielddesc:EFUSEFLASHTPUW]{EFUSE\_FLASH\_TPUW} & 4 & N & 18 & Represents flash waiting time after power-up. \\    
    \hline    
    \hyperref[fielddesc:EFUSEDISDOWNLOADMODE]{EFUSE\_DIS\_DOWNLOAD\_MODE} & 1 & N & 18 & Represents whether all download modes are disabled.\\       
    \hline
    \hyperref[fielddesc:EFUSEDISLEGACYSPIBOOT]{EFUSE\_DIS\_LEGACY\_SPI\_BOOT} & 1 & N & 18 & Represents whether Legacy SPI is disabled. \\
    \hline
    \hyperref[fielddesc:EFUSEDISUSBPRINT]{EFUSE\_DIS\_USB\_PRINT} & 1 & N & 18 & Represents whether USB printing is disabled. \\
    \hline  
    \hyperref[fielddesc:EFUSEFLASHECCMODE]{EFUSE\_FLASH\_ECC\_MODE} & 1 & N & 18 & Represents the flash ECC mode in ROM. \\
    \hline
    \hyperref[fielddesc:EFUSEDISUSBSERIALJTAGDOWNLOADMODE]{EFUSE\_DIS\_USB\_SERIAL\_JTAG\_DOWNLOAD\_MODE} & 1 & N & 18 & Represents whether download through USB-Serial-JTAG is disabled. \\
    \hline 
    \hyperref[fielddesc:EFUSEENABLESECURITYDOWNLOAD]{EFUSE\_ENABLE\_SECURITY\_
    DOWNLOAD} & 1 & N & 18 & Represents whether secure UART download mode is enabled. \\
    \hline
    \hyperref[fielddesc:EFUSEUARTPRINTCONTROL]{EFUSE\_UART\_PRINT\_CONTROL} & 2 & N & 18 & Represents the default UART boot message output mode. \\
    \hline
    \hyperref[fielddesc:EFUSEPINPOWERSELECTION]{EFUSE\_PIN\_POWER\_SELECTION} & 1 & N & 18 & Represents the power supply for GPIO33 \verb+~+ GPIO37, GPIO47, and GPIO48. \\
    \hline
    \hyperref[fielddesc:EFUSEFLASHTYPE]{EFUSE\_FLASH\_TYPE} & 1 & N & 18 & Represents the maximum data lines of SPI flash. \\
    \hline
    \hyperref[fielddesc:EFUSEFLASHPAGESIZE]{EFUSE\_FLASH\_PAGE\_SIZE} & 2 & N & 18 & Represents the page size of flash.\\
    \hline
    \hyperref[fielddesc:EFUSEFLASHECCEN]{EFUSE\_FLASH\_ECC\_EN} & 1 & N & 18 & Represents whether ECC for flash boot is enabled. \\
    \hline
    \hyperref[fielddesc:EFUSEFORCESENDRESUME]{EFUSE\_FORCE\_SEND\_RESUME} & 1 & N & 18 & Represents whether to force ROM code to send a resume command during SPI boot. \\  
    \hline
    \hyperref[fielddesc:EFUSESECUREVERSION]{EFUSE\_SECURE\_VERSION} & 16 & N & 18 & Represents IDF secure version. \\\hline
    \hyperref[fielddesc:EFUSEDISUSBOTGDOWNLOADMODE]{EFUSE\_DIS\_USB\_OTG\_DOWNLOAD\_MODE} & 1 & N & 19 & Represents whether download through USB-OTG is disabled. \\\hline
\end{longtable}


\end{landscape}
\restoregeometry
\clearpage


Table \ref{tab:efuse-key-purpose} lists all key purpose and their values. Setting the eFuse parameter EFUSE\_KEY\_PURPOSE\_\textcolor{red}{\it{n}} declares the purpose of KEY\textcolor{red}{\it{n}} (\textcolor{red}{\it{n}}: 0 \verb+~+ 5).

\begin{longtable}[c]{ | C{1.5cm} | p{10cm} | }
\caption{Secure Key Purpose Values}
\label{tab:efuse-key-purpose}
    \\\hline
    \rowcolor{lightgray}
        \textbf{Key Purpose Values} &  \textbf{Purposes} \\
    \endfirsthead
    
    \multicolumn{2}{c}{\bfseries \tablename\ \thetable{} -- cont'd from previous page}\\\hline
    
    \rowcolor{lightgray}
        \textbf{Key Purpose Values} &  \textbf{Purposes} \\
    \endhead
    \multicolumn{2}{r}{\textbf{Cont'd on next page}} \\
    \endfoot
    \endlastfoot
    \hline
0 & User purposes\\\hline
1 &  Reserved \\\hline
2 & XTS\_AES\_256\_KEY\_1 (flash/SRAM encryption and decryption) \\\hline
3 & XTS\_AES\_256\_KEY\_2 (flash/SRAM encryption and decryption) \\ \hline
4 & XTS\_AES\_128\_KEY (flash/SRAM encryption and decryption) \\\hline
5 & HMAC Downstream mode \\\hline
6 & JTAG in HMAC Downstream mode \\\hline
7 & RSA Digital Signature Peripheral (RSA\_DS) in HMAC Downstream mode \\\hline
8 & HMAC Upstream mode \\\hline
9 & SECURE\_BOOT\_DIGEST0 (secure boot key digest) \\\hline
10 & SECURE\_BOOT\_DIGEST1 (secure boot key digest)\\\hline
11 & SECURE\_BOOT\_DIGEST2 (secure boot key digest) \\\hline
\end{longtable}

Table  \ref{tab:efuse-block-1-10-para} provides the details of parameters in BLOCK1 \verb+~+ BLOCK10.

\clearpage
\begin{landscape}

\input{\modulefiles/tables/07-eFuse-block-1-10-para__EN}

\end{landscape}
\restoregeometry
\clearpage


Among these blocks, BLOCK4 \verb+~+ 9 store KEY0 \verb+~+ 5, respectively. Up to six 256-bit keys can be written into eFuse. Whenever a key is written, its purpose value should also be written (see table \ref{tab:efuse-key-purpose}). For example, when a key for the JTAG function in HMAC Downstream mode is written to KEY\textcolor{red}{3} (i.e., BLOCK7), its key purpose value 6 should also be written to EFUSE\_KEY\_PURPOSE\_\textcolor{red}{3}.

BLOCK1 \verb+~+ BLOCK10 use the RS coding scheme, so there are some restrictions on writing to these parameters. For more detailed information, please refer to Section \ref{sec:efuse-data-storage} and \ref{sec:efuse-para-software-programming}.

\subsubsection{\hyperref[fielddesc:EFUSEWRDIS]{EFUSE\_WR\_DIS}}
\label{sec:efuse-wr-dis}

Parameter \hyperref[fielddesc:EFUSEWRDIS]{EFUSE\_WR\_DIS} determines whether individual eFuse parameters are write-protected. After \hyperref[fielddesc:EFUSEWRDIS]{EFUSE\_WR\_DIS} has been programmed, execute an eFuse read operation to let the new values take effect (see Section \ref{sec:efuse-software-read-para}).

Column ``Write Protection by \hyperref[fielddesc:EFUSEWRDIS]{EFUSE\_WR\_DIS} Bit Number" in Table  \ref{tab:efuse-block0-para} and Table \ref{tab:efuse-block-1-10-para} list the specific bits in \hyperref[fielddesc:EFUSEWRDIS]{EFUSE\_WR\_DIS} that disable writing.

When the write protection bit of a parameter is set to 0, it means that this parameter is not write-protected and can be programmed.

Setting the write protection bit of a parameter to 1 enables write-protection for it and none of its bits can be modified afterwards. Non-programmed bits always remain 0 while programmed bits always remain 1.

\subsubsection{\hyperref[fielddesc:EFUSERDDIS]{EFUSE\_RD\_DIS}}
\label{sec:efuse-rd-dis}

Only the eFuse blocks BLOCK4 \verb+~+ BLOCK10 can be individually read protected to prevent any access from outside the chip, as shown in column ``Read Protection by \hyperref[fielddesc:EFUSERDDIS]{EFUSE\_RD\_DIS} Bit Number" of Table  \ref{tab:efuse-block-1-10-para}. After \hyperref[fielddesc:EFUSERDDIS]{EFUSE\_RD\_DIS} has been programmed, execute an eFuse read operation to let the new values take effect (see Section \ref{sec:efuse-software-read-para}).

If a bit in \hyperref[fielddesc:EFUSERDDIS]{EFUSE\_RD\_DIS} is 0, then the eFuse block can be read by users; if a bit in \hyperref[fielddesc:EFUSERDDIS]{EFUSE\_RD\_DIS} is 1, then the parameter controlled by this bit is user read protected.

Other parameters that are not in BLOCK4 \verb+~+ BLOCK10 can always be read by users.

When BLOCK4 \verb+~+ BLOCK10 are set to be read-protected, the data in these blocks are not readable by users, but they can still be used internally by hardware cryptography modules, if the EFUSE\_KEY\_PURPOSE\_\textcolor{red}{\it{n}} bit is set accordingly.

\subsubsection{Data Storage}
\label{sec:efuse-data-storage}

According to the different types of eFuse bits, eFuse controller use two hardware encoding schemes to protect eFuse bits from corruption.

All BLOCK0 parameters except for \hyperref[fielddesc:EFUSEWRDIS]{EFUSE\_WR\_DIS} are stored with four backups, meaning each bit is stored four times. This scheme is transparent to the user. This encoding scheme is invisible for users.

BLOCK1 \verb+~+ BLOCK10 store key data and some parameters and use RS (44, 32) coding scheme that supports up to 6 bytes of automatic error correction. The primitive polynomial of RS (44, 32) is \(p(x)=x^8+x^4+x^3+x^2+1\).

\begin{figure}[ht!]
    \centering
    \fbox{\includegraphics[width=1.0\textwidth]{\modulefiles/figures/07-eFuse-shift-reg-cir-1.eps}}
    \caption{Shift Register Circuit (output of first 32 bytes)}
    \label{fig:efuse-lfsr-circuit-1}
\end{figure}

\begin{figure}[ht!]
    \centering
    \fbox{\includegraphics[width=1.0\textwidth]{\modulefiles/figures/07-eFuse-shift-reg-cir-2.eps}}
    \caption{Shift Register Circuit (output of last 12 bytes)}
    \label{fig:efuse-lfsr-circuit-2}
\end{figure}


The shift register circuit shown in Figure \ref{fig:efuse-lfsr-circuit-1} and \ref{fig:efuse-lfsr-circuit-2} processes 32 data bytes using RS (44, 32). This coding scheme encodes 32 bytes of  data into 44 bytes:

\begin{itemize}
    \item Bytes [0:31] are the data bytes itself
    \item Bytes [32:43] are the encoded parity bytes stored in 8-bit flip-flops DFF1, DFF2, ..., DFF12 (gf\_mul\_n, where n is an integer, is the result of multiplying a byte of data ...)
\end{itemize}

After that, the hardware burns into eFuse the 44-byte codeword consisting of the data bytes followed by the parity bytes.

When the eFuse block is read back, the eFuse controller automatically decodes the codeword  and applies error correction if needed.

Because the RS check codes are generated on the entire 256-bit eFuse block, each block can only be written once.

\subsection{Programming of Parameters}
\label{sec:efuse-para-software-programming}

The eFuse controller can only program eFuse parameters of one block at a time. BLOCK0 \verb+~+ BLOCK10 share the same address range to store the parameters to be programmed. Configure parameter \hyperref[fielddesc:EFUSEBLKNUM]{EFUSE\_BLK\_NUM} to indicate which block should be programmed.

Before programming, make sure the eFuse programming voltage VDDQ is configured correctly as described in Section \ref{sec:efuse-vddq-timing}.

{\bfseries Programming BLOCK0}
\label{sec:block0_hw_bits}

When \hyperref[fielddesc:EFUSEBLKNUM]{EFUSE\_BLK\_NUM} is set to 0, BLOCK0 will be programmed. Register \hyperref[regdesc:EFUSEPGMDATA0REG]{EFUSE\_PGM\_DATA0\_REG} stores \hyperref[fielddesc:EFUSEWRDIS]{EFUSE\_WR\_DIS}. Registers \hyperref[regdesc:EFUSEPGMDATA1REG]{EFUSE\_PGM\_DATA1\_REG} \verb+~+ \hyperref[regdesc:EFUSEPGMDATA5REG]{EFUSE\_PGM\_DATA5\_REG} store the information of parameters to be programmed. Note that 25 bits are readable but useless to users and must always be set to 0 in the programming registers. The specific bits are:

\begin{itemize}
    \item \hyperref[regdesc:EFUSEPGMDATA1REG]{EFUSE\_PGM\_DATA1\_REG}[27:31]
    \item \hyperref[regdesc:EFUSEPGMDATA1REG]{EFUSE\_PGM\_DATA1\_REG}[21:24]
    \item \hyperref[regdesc:EFUSEPGMDATA2REG]{EFUSE\_PGM\_DATA2\_REG}[7:15]
    \item \hyperref[regdesc:EFUSEPGMDATA2REG]{EFUSE\_PGM\_DATA2\_REG}[0:3]
    \item \hyperref[regdesc:EFUSEPGMDATA3REG]{EFUSE\_PGM\_DATA3\_REG}[26:27]
    \item \hyperref[regdesc:EFUSEPGMDATA4REG]{EFUSE\_PGM\_DATA4\_REG}[30]
\end{itemize}

Data in registers \hyperref[regdesc:EFUSEPGMDATA6REG]{EFUSE\_PGM\_DATA6\_REG} \verb+~+ \hyperref[regdesc:EFUSEPGMDATA7REG]{EFUSE\_PGM\_DATA7\_REG} and \hyperref[regdesc:EFUSEPGMCHECKVALUE0REG]{EFUSE\_PGM\_CHECK\_VALUE0\_REG} \verb+~+ \hyperref[regdesc:EFUSEPGMCHECKVALUE2REG]{EFUSE\_PGM\_CHECK\_VALUE2\_REG} are ignored when programming BLOCK0.

{\bfseries Programming BLOCK1}

When \hyperref[fielddesc:EFUSEBLKNUM]{EFUSE\_BLK\_NUM} is set to 1, registers \hyperref[regdesc:EFUSEPGMDATA0REG]{EFUSE\_PGM\_DATA0\_REG} \verb+~+ \hyperref[regdesc:EFUSEPGMDATA5REG]{EFUSE\_PGM\_DATA5\_REG} store the BLOCK1 parameters to be programmed. Registers \hyperref[regdesc:EFUSEPGMCHECKVALUE0REG]{EFUSE\_PGM\_CHECK\_VALUE0\_REG} \verb+~+ \hyperref[regdesc:EFUSEPGMCHECKVALUE2REG]{EFUSE\_PGM\_CHECK\_VALUE2\_REG} store the corresponding RS check codes. Data in registers \hyperref[regdesc:EFUSEPGMDATA6REG]{EFUSE\_PGM\_DATA6\_REG} \verb+~+ \hyperref[regdesc:EFUSEPGMDATA7REG]{EFUSE\_PGM\_DATA7\_REG} is ignored when programming BLOCK1, and the RS check codes will be calculated with these bits all treated as 0.

{\bfseries Programming BLOCK2 \verb+~+ 10}

When \hyperref[fielddesc:EFUSEBLKNUM]{EFUSE\_BLK\_NUM} is set to 2 \verb+~+ 10, registers \hyperref[regdesc:EFUSEPGMDATA0REG]{EFUSE\_PGM\_DATA0\_REG} \verb+~+ \hyperref[regdesc:EFUSEPGMDATA7REG]{EFUSE\_PGM\_DATA7\_REG} store the parameters to be programmed to this block. Registers\\ \hyperref[regdesc:EFUSEPGMCHECKVALUE0REG]{EFUSE\_PGM\_CHECK\_VALUE0\_REG} \verb+~+ \hyperref[regdesc:EFUSEPGMCHECKVALUE2REG]{EFUSE\_PGM\_CHECK\_VALUE2\_REG} store the corresponding RS check codes.

{\bfseries Programming process}

The process of programming parameters is as follows:

\begin{enumerate}
\item Write the block number to \hyperref[fielddesc:EFUSEBLKNUM]{EFUSE\_BLK\_NUM} to determine the block to be programmed.

\item Write parameters to be programmed to registers \hyperref[regdesc:EFUSEPGMDATA0REG]{EFUSE\_PGM\_DATA0\_REG} \verb+~+ \hyperref[regdesc:EFUSEPGMDATA7REG]{EFUSE\_PGM\_DATA7\_REG} and the corresponding checksum values to\\ \hyperref[regdesc:EFUSEPGMCHECKVALUE0REG]{EFUSE\_PGM\_CHECK\_VALUE0\_REG} \verb+~+ \hyperref[regdesc:EFUSEPGMCHECKVALUE2REG]{EFUSE\_PGM\_CHECK\_VALUE2\_REG}.

\item Configure the field \hyperref[fielddesc:EFUSEOPCODE]{EFUSE\_OP\_CODE} of register \hyperref[regdesc:EFUSECONFREG]{EFUSE\_CONF\_REG} to 0x{}5A5A.

\item Configure the field \hyperref[fielddesc:EFUSEPGMCMD]{EFUSE\_PGM\_CMD} of register \hyperref[regdesc:EFUSECMDREG]{EFUSE\_CMD\_REG} to 1.

\item Poll register \hyperref[regdesc:EFUSECMDREG]{EFUSE\_CMD\_REG} until it is 0x{}0, or wait for a PGM\_DONE interrupt. For more information on how to identify a PGM\_DONE interrupt, please see the end of Section \ref{sec:efuse-software-read-para}.

\item In order to avoid programming content leakage, please clear the parameters in \hyperref[regdesc:EFUSEPGMDATA0REG]{EFUSE\_PGM\_DATA0\_REG} \verb+~+ \hyperref[regdesc:EFUSEPGMDATA7REG]{EFUSE\_PGM\_DATA7\_REG} and \\ \hyperref[regdesc:EFUSEPGMCHECKVALUE0REG]{EFUSE\_PGM\_CHECK\_VALUE0\_REG} \verb+~+ \hyperref[regdesc:EFUSEPGMCHECKVALUE2REG]{EFUSE\_PGM\_CHECK\_VALUE2\_REG}.

\item Trigger an eFuse read operation (see Section \ref{sec:efuse-software-read-para}) to update eFuse registers with the new values.
\item Check error record registers. If the values read in error record registers are not 0, the programming process should be performed again following above steps 1 \verb+~+ 7. Please check the following error record registers for different eFuse blocks:
\begin{itemize}
    \item BLOCK0: \hyperref[regdesc:EFUSERDREPEATERR0REG]{EFUSE\_RD\_REPEAT\_ERR0\_REG} \verb+~+ \hyperref[regdesc:EFUSERDREPEATERR4REG]{EFUSE\_RD\_REPEAT\_ERR4\_REG}
    \item BLOCK1: \hyperref[regdesc:EFUSERDRSERR0REG]{EFUSE\_RD\_RS\_ERR0\_REG}[2:0], \hyperref[regdesc:EFUSERDRSERR0REG]{EFUSE\_RD\_RS\_ERR0\_REG}[7]
    \item BLOCK2: \hyperref[regdesc:EFUSERDRSERR0REG]{EFUSE\_RD\_RS\_ERR0\_REG}[6:4], \hyperref[regdesc:EFUSERDRSERR0REG]{EFUSE\_RD\_RS\_ERR0\_REG}[11]
    \item BLOCK3: \hyperref[regdesc:EFUSERDRSERR0REG]{EFUSE\_RD\_RS\_ERR0\_REG}[10:8], \hyperref[regdesc:EFUSERDRSERR0REG]{EFUSE\_RD\_RS\_ERR0\_REG}[15]
    \item BLOCK4: \hyperref[regdesc:EFUSERDRSERR0REG]{EFUSE\_RD\_RS\_ERR0\_REG}[14:12], \hyperref[regdesc:EFUSERDRSERR0REG]{EFUSE\_RD\_RS\_ERR0\_REG}[19]
    \item BLOCK5: \hyperref[regdesc:EFUSERDRSERR0REG]{EFUSE\_RD\_RS\_ERR0\_REG}[18:16], \hyperref[regdesc:EFUSERDRSERR0REG]{EFUSE\_RD\_RS\_ERR0\_REG}[23]
    \item BLOCK6: \hyperref[regdesc:EFUSERDRSERR0REG]{EFUSE\_RD\_RS\_ERR0\_REG}[22:20], \hyperref[regdesc:EFUSERDRSERR0REG]{EFUSE\_RD\_RS\_ERR0\_REG}[27]
    \item BLOCK7: \hyperref[regdesc:EFUSERDRSERR0REG]{EFUSE\_RD\_RS\_ERR0\_REG}[26:24], \hyperref[regdesc:EFUSERDRSERR0REG]{EFUSE\_RD\_RS\_ERR0\_REG}[31]
    \item BLOCK8: \hyperref[regdesc:EFUSERDRSERR0REG]{EFUSE\_RD\_RS\_ERR0\_REG}[30:28], \hyperref[regdesc:EFUSERDRSERR1REG]{EFUSE\_RD\_RS\_ERR1\_REG}[3]
    \item BLOCK9: \hyperref[regdesc:EFUSERDRSERR1REG]{EFUSE\_RD\_RS\_ERR1\_REG}[2:0], \hyperref[regdesc:EFUSERDRSERR1REG]{EFUSE\_RD\_RS\_ERR1\_REG}[2:0][7]
    \item BLOCK10: \hyperref[regdesc:EFUSERDRSERR1REG]{EFUSE\_RD\_RS\_ERR1\_REG}[2:0][6:4]
\end{itemize}
\end{enumerate}

{\bfseries Limitations}

In BLOCK0, each bit can be programmed separately. However, we recommend to minimize programming cycles and program all the bits of a parameter in one programming action. In addition, after all parameters controlled by a certain bit of \hyperref[fielddesc:EFUSEWRDIS]{EFUSE\_WR\_DIS} are programmed, that bit should be immediately programmed. The programming of parameters controlled by a certain bit of \hyperref[fielddesc:EFUSEWRDIS]{EFUSE\_WR\_DIS},
and the programming of the bit itself can even be completed at the same time. Repeated programming of already programmed bits is strictly forbidden, otherwise, programming errors will occur.

BLOCK1 cannot be programmed by users as it has been programmed at manufacturing.

BLOCK2 \verb+~+ 10 can only be programmed once. Repeated programming is not allowed.

\subsection{User Read of Parameters}
\label{sec:efuse-software-read-para}

Users cannot read eFuse bits directly. The eFuse Controller hardware reads all eFuse bits and stores the results to their corresponding registers in its memory space. Then, users can read eFuse bits by reading the registers that start with EFUSE\_RD\_. Details are provided in Table \ref{tab:efuse-softw-read-para}.

\input{\modulefiles/tables/07-eFuse-software-read-para__EN}

{\bfseries Updating eFuse read registers}
\label{sec:efuse-update-reg}

The eFuse Controller reads internal eFuses to update corresponding registers. This read operation happens on system reset and can also be triggered manually by users as needed (e.g., if new eFuse values have been programmed). The process of triggering a read operation by users is as follows:

\begin{enumerate}
    \item Configure the field \hyperref[fielddesc:EFUSEOPCODE]{EFUSE\_OP\_CODE} in register \hyperref[regdesc:EFUSECONFREG]{EFUSE\_CONF\_REG} to 0x{}5AA5.
    \item Configure the field \hyperref[fielddesc:EFUSEREADCMD]{EFUSE\_READ\_CMD} in register \hyperref[regdesc:EFUSECMDREG]{EFUSE\_CMD\_REG} to 1.
    \item Poll register \hyperref[regdesc:EFUSECMDREG]{EFUSE\_CMD\_REG} until it is 0x{}0, or wait for a READ\_DONE interrupt. Information on how to identify a READ\_DONE interrupt is provided below in this section.
    \item users reads the values of each parameter from memory.
\end{enumerate}
The eFuse read registers will hold all values until the next read operation.

{\bfseries Error detection}

Error record registers allow users to detect if there are any inconsistencies in the stored backup eFuse parameters.

Registers EFUSE\_RD\_REPEAT\_ERR0 \verb+~+ 3\_REG indicate if there are any errors of programmed parameters (except for \hyperref[fielddesc:EFUSEWRDIS]{EFUSE\_WR\_DIS}) in BLOCK0 (value 1 indicates an error is detected, and the bit becomes invalid; value 0 indicates no error).

Registers EFUSE\_RD\_RS\_ERR0 \verb+~+ 1\_REG store the number of corrected bytes as well as the result of RS decoding during eFuse reading BLOCK1 \verb+~+ BLOCK10.

The values of above registers will be updated every time after the eFuse read registers have been updated.

{\bfseries Identifying the completion of a program/read operation}

    The methods to identify the completion of a program/read operation are described below. Please note that bit 1 corresponds to a program operation, and bit 0 corresponds to a read operation.
    \begin{itemize}
    \item Method one:
        \begin{enumerate}
            \item Poll bit 1/0 in register \hyperref[regdesc:EFUSEINTRAWREG]{EFUSE\_INT\_RAW\_REG} until it becomes 1, which represents the completion of a program/read operation.
        \end{enumerate}
    \item Method two:
        \begin{enumerate}
            \item Set bit 1/0 in register \hyperref[regdesc:EFUSEINTENAREG]{EFUSE\_INT\_ENA\_REG} to 1 to enable the eFuse Controller to post a PGM\_DONE or READ\_DONE interrupt.
            \item Configure the Interrupt Matrix to enable the CPU to respond to eFuse interrupt signals, see Chapter \ref{mod:intmtrx} \textit{\nameref{mod:intmtrx}}.
            \item Wait for the PGM/READ\_DONE interrupt.
            \item Set bit 1/0 in register \hyperref[regdesc:EFUSEINTCLRREG]{EFUSE\_INT\_CLR\_REG} to 1 to clear the PGM/READ\_DONE interrupt.
    \end{enumerate}
\end{itemize}

{\bfseries Note}


When eFuse controller updating its registers, it will use \hyperref[regdesc:EFUSEPGMDATA0REG]{EFUSE\_PGM\_DATA\textcolor{red}{\it{n}}\_REG} (n=0, 1, .., 7) again to store data. So please do not write important data into these registers before this updating process initiated. \\
During the chip boot process, eFuse controller will update eFuse data into registers which can be accessed by users automatically. You can get programmed eFuse data by reading corresponding registers. Thus, it is no need to update eFuse read registers in such case.

\subsection{eFuse VDDQ Timing}
\label{sec:efuse-vddq-timing}

The eFuse Controller operates with 20 MHz, one cycle is 50 ns, and its programming voltage VDDQ should be configured as follows:

\begin{itemize}
    \item \hyperref[fielddesc:EFUSEDACNUM]{EFUSE\_DAC\_NUM} (store the rising period of VDDQ): The default value of VDDQ is 2.5 V and the voltage increases by 0.01 V in each clock cycle. Thus, the default value of this parameter is 255;
    \item \hyperref[fielddesc:EFUSEDACCLKDIV]{EFUSE\_DAC\_CLK\_DIV} (the clock divisor of VDDQ): The clock period to program VDDQ should be larger than 1 $\mu$s;
    \item \hyperref[fielddesc:EFUSEPWRONNUM]{EFUSE\_PWR\_ON\_NUM} (the power-up time for VDDQ): The programming voltage should be stabilized after this time, which means the value of this parameter should be configured to exceed the value of \hyperref[fielddesc:EFUSEDACCLKDIV]{EFUSE\_DAC\_CLK\_DIV} multiply by \hyperref[fielddesc:EFUSEDACNUM]{EFUSE\_DAC\_NUM};
    \item \hyperref[fielddesc:EFUSEPWROFFNUM]{EFUSE\_PWR\_OFF\_NUM} (the power-out time for VDDQ): The value of this parameter should be larger than 10 $\mu$s.
\end{itemize}

\input{\modulefiles/tables/07-eFuse-vddq-tim-config__EN}

\subsection{The Use of Parameters by Hardware Modules}

Some hardware modules are directly connected to the eFuse peripheral in order to use the parameters listed in Table \ref{tab:efuse-block0-para} and Table \ref{tab:efuse-block-1-10-para}, specifically those marked with ``Y" in columns ``Accessible by Hardware". Users cannot intervene in this process.

\subsection{Interrupts}
\begin{itemize}
\item PGM\_DONE interrupt: Triggered when eFuse programming has finished. Set \hyperref[fielddesc:EFUSEPGMDONEINTENA]{EFUSE\_PGM\_DONE\_INT\_ENA} to enable this interrupt;
\item READ\_DONE interrupt: Triggered when eFuse reading has finished. Set \hyperref[fielddesc:EFUSEREADDONEINTENA]{EFUSE\_READ\_DONE\_INT\_ENA} to enable this interrupt.
\end{itemize}


% >>> If your module has no registers, delete sample text
%       for Sections Register Summary and Registers below <<<

%%%%%%%%%%%%%%%%%%%%%%%%
%%% Register Summary %%%
%%%%%%%%%%%%%%%%%%%%%%%%

% >>> Update the sample text below and insert
%     the info on register summary into the subfile <<<

\clearpage
\hypertarget{efuse-reg-summ}{}
\section{Register Summary}
\label{sec:efuse-reg-summ}

The addresses in this section are relative to eFuse Controller base address provided in Table \ref{tab:sysmem-base-address} in Chapter \ref{mod:sysmem} \textit{\nameref{mod:sysmem}}.

The abbreviations given in Column \textbf{Access} are explained in Section \hyperref[glossary-access-types]{\textit{Access Types for Registers}}.

\subfile{\modulefiles/07-eFuse-reg-summ__EN}


%%%%%%%%%%%%%%%%%
%%% Registers %%%
%%%%%%%%%%%%%%%%%

% >>> Update the sample text below and insert
%      the info on registers into the subfile <<<

\clearpage
\section{Registers}

The addresses in this section are relative to eFuse Controller base address provided in Table \ref{tab:sysmem-base-address} in Chapter \ref{mod:sysmem} \textit{\nameref{mod:sysmem}}.

\subfile{\modulefiles/07-eFuse-reg__EN}


\end{document}
